\section{Regional states under relabelling}
\label{sec:area_law_regional_relabelling}

\begin{definition}[Relabelling physical configurations]
    \label{def:area_law_regional_relabelling}
    \lean{TNLean.PEPS.AreaLaw.configurationSiteEquiv,
        TNLean.PEPS.AreaLaw.regionConfigurationEquiv,
        TNLean.PEPS.AreaLaw.relabelState}
    \leanok
    \uses{def:area_law_domain,def:area_law_ground_entropy}
    Let $e:\Lambda\to\Lambda'$ be a bijection of finite domains, and let
    $A\subseteq\Lambda$, $A'\subseteq\Lambda'$ satisfy $e(A)=A'$.
    It induces bijections $e_A:A\to A'$ and
    $e_B:\Lambda\setminus A\to\Lambda'\setminus A'$.
    For any physical vector $\Omega$, its relabelled coefficients are
    \begin{align}
        \Omega'(\sigma') & =\Omega(\sigma'\circ e).
    \end{align}
    Regional configurations transform by the corresponding restricted bijections.
    The relabelling concerns tensor factors; preservation of lattice adjacency
    is unnecessary for this definition.
\end{definition}

\begin{theorem}[Covariance of reductions and invariance of entropy]
    \label{thm:area_law_regional_relabelling}
    \lean{TNLean.PEPS.AreaLaw.reducedState_relabel,
        TNLean.PEPS.AreaLaw.regionalEntropy_relabel}
    \leanok
    \uses{def:area_law_regional_relabelling}
    For every vector and region-preserving site bijection as above,
    \begin{align}
        \|\Omega'\|     & =\|\Omega\|,                                     \\
        \rho_{A'}^{\Omega'}(\sigma',\tau')
                        & =\rho_A^\Omega(\sigma'\circ e_A,\tau'\circ e_A), \\
        S_{\Omega'}(A') & =S_\Omega(A).
    \end{align}
    These identities hold without normalization or spectral assumptions,
    including empty regions and complements. They express the finite tensor-factor
    conventions of \cite[Section~2]{OpenAI2026AreaLaw} in the induced-domain coordinates.
\end{theorem}
\begin{proof}\leanok
    Relabelling permutes the Euclidean configuration coordinates, preserving the
    norm. The site bijection preserves the region and its complement. Splitting a
    configuration before or after relabelling therefore gives the same pair of
    regional configurations. The partial trace changes coordinates on both
    factors, while its complementary sum is unchanged by the complementary
    bijection. This gives the reduction formula. Reindexing the regional matrix
    preserves its eigenvalues and hence its von Neumann entropy.
\end{proof}
